\documentclass[a4paper,12pt]{article}

% === Кодировки и шрифты для русского языка ===
\usepackage[T2A]{fontenc}
\usepackage[utf8]{inputenc}
\usepackage[russian]{babel}
\usepackage{indentfirst}

% === Математика и графика ===
\usepackage{amsmath, amssymb, amsfonts, mathtools, bm}
\usepackage{graphicx, float, booktabs, multirow}

% === Форматирование ===
\usepackage[left=2.5cm, right=1.5cm, top=2cm, bottom=2cm]{geometry}
\usepackage{setspace}
\onehalfspacing

% === Ссылки и библиография ===
\usepackage{cite, url, hyperref}
\hypersetup{colorlinks=true, linkcolor=black, citecolor=black, urlcolor=blue}
\usepackage{newunicodechar}
\newunicodechar{́}{\'}

% === Нумерация формул по разделам ===
\numberwithin{equation}{section}

% === Заголовок и авторы ===
\title{\textbf{Частотное разнесение в OFDM-системах через управляемый алиасинг спектральных образов}}

\author{
    Рычков~Е.~Н.$^{1}$ \\
    $^{1}$Сибирский Федеральный университет, \texttt{meurchmusic@gmail.com} \\
}
\date{}

\begin{document}

\maketitle

% === Аннотация ===
\begin{abstract}
\noindent \textbf{Аннотация} --- Предложена архитектура частотного разнесения для систем с ортогональным частотным разделением (OFDM), использующая спектральные образы в зонах Найквиста без повышения частоты дискретизации или применения дополнительных гетеродинов. Применение двух аналого-цифровых преобразователей (АЦП) с временным чередованием каналов со сдвигом тактовой фазы на полпериода ($\Delta t = T_s/2$) обеспечивает детерминированный фазовый сдвиг $\pi$ между копиями сигнала из соседних зон Найквиста, что позволяет их обратимое цифровое разделение посредством преобразования Адамара. После разделения когерентное комбинирование независимых ветвей методом максимального отношения (MRC) обеспечивает выигрыш частотного разнесения в условиях рэлеевских замираний. Работа в широкополосном режиме обеспечивается поподнесущей фазовой коррекцией, компенсирующей частотно-зависимый набег фазы от сдвига тактов. Результаты Монте-Карло моделирования с учётом типовых неидеальностей АЦП подтверждают порядок частотного разнесения $d = 1{,}896 \pm 0{,}020$. В области сверхнизких вероятностей ошибки ($\text{BER}=10^{-5}$) преимущество метода сохраняется, обеспечивая снижение требуемого отношения сигнал/шум на $\sim$12 дБ по сравнению с современной типовой архитектурой уже при M=2 (основной базовый спектр и первая следующая зона Найквиста). Подробности эксперимента, включая исходный код, параметры симуляции и скрипты для воспроизведения результатов, приведены в репозитории~\cite{github}. Метод применим в сценариях с ограниченной мощностью передатчика и относительно неограниченным спектром: системы Интернета вещей, подводная и космическая связь.

\medskip
\noindent \textbf{Ключевые слова}: частотное разнесение, спектральные образы, зоны Найквиста, интерливинговые АЦП, управляемый алиасинг, матрица Адамара, когерентное комбинирование, помехоустойчивость.
\end{abstract}

% === 1. Введение ===
\section{Введение}
\label{sec:intro}

Генерация широкополосных сигналов в системах связи пятого и шестого поколений (40--70~ГГц, полоса сигнала 1--2~ГГц) является актуальной задачей, решение которой традиционными методами требует либо высокоскоростных цифро-аналоговых преобразователей (ЦАП) с тактовой частотой, превышающей удвоенную максимальную частоту сигнала (согласно теореме Котельникова), либо применения сложных многоканальных схем с гетеродинами. Оба подхода существенно увеличивают стоимость, энергопотребление и сложность оборудования.

Известно, что на выходе ЦАП при формировании аналогового сигнала возникают периодические спектральные копии (образы) в зонах Найквиста с номерами $N = 1, 2, 3\dots$, отделённые друг от друга на величину частоты дискретизации $f_s$ \cite{oppenheim} и затухающие по амплитуде от зоны к зоне ввиду импульсной природы в ЦАП при невозможности использования идеальной дельта-функции. 

В традиционных системах связи эти образы рассматриваются как паразитные гармоники и подавляются фильтрацией нижних частот, хотя активно исследуется направление (ЦАП AD9119/9129, EV12DS130А) возможности работы в N-зонах Найквиста. Также в анализаторах спектра можно обнаружить технику поочередного семплирования двумя и более АЦП для увеличения полосы сигнала. В данной работе предлагается использовать спектральные образы не как артефакты, а как дополнительные частотно-разнесённые копии полезного сигнала для повышения помехоустойчивости.

\subsection*{Факторы, препятствовавшие развитию подхода}
Несмотря на фундаментальную известность зон Найквиста, целенаправленное использование спектральных образов для частотного разнесения не исследовалось в литературе по следующим причинам:
\begin{itemize}
    \item \textbf{Терминологическая фрагментация}: метод лежит на стыке концепций \textit{bandpass sampling} \cite{vaughan,rajasekaran}, \textit{Nyquist folding} \cite{fudge,murray}, \textit{multi-Nyquist zone}, \textit{image-band exploitation}, \textit{Direct RF sampling OFDM}, \textit{fractional sampling diversity} \cite{gadawe}, \textit{Sub-Nyquist OFDM} \cite{han}, что затрудняет его идентификацию и систематизацию;
    \item \textbf{Фокус на расширение полосы}: большинство предшествующих работ использовали высшие зоны Найквиста исключительно для синтеза широкой полосы сигнала \cite{vaughan,fudge}, а не для создания избыточности с целью разнесения;
    \item \textbf{Проблема фазовых соотношений}: до развития технологий интерливинговых АЦП и методов цифровой калибровки фазовые соотношения между спектральными копиями считались неконтролируемыми и непригодными для детерминированной обработки.
\end{itemize}

\subsection*{Основные положения метода}
\begin{enumerate}
    \item ЦАП при генерации OFDM-сигнала автоматически создаёт спектральные копии в зонах Найквиста $N=1,2,3\dots$;
    \item При оцифровке двумя АЦП со сдвигом тактовой фазы $\Delta t = T_s/2$ копии из соседних зон приобретают детерминированный фазовый сдвиг $\pi$;
    \item Это позволяет обратимо разделить копии матрицей Адамара и когерентно сложить их мощности через MRC, повышая помехоустойчивость без искажений и без повышения тактовой частоты.
\end{enumerate}

\textbf{Основной вклад работы}:
\begin{enumerate}
    \item Предложена архитектура приёмника с двумя интерливинговыми АЦП, сдвиг тактовой фазы между которыми на $\Delta t = T_s/2$ создаёт детерминированный фазовый сдвиг $\pi$ между спектральными копиями из соседних зон Найквиста, что позволяет их обратимое цифровое разделение.
    \item Разработан алгоритм поподнесущей фазовой коррекции для обеспечения работы метода в широкополосном режиме OFDM (компенсация частотно-зависимого набега фазы).
    \item Проведена верификация метода методом Монте-Карло с учётом типовых неидеальностей АЦП, подтвердившая достижение порядка частотного разнесения $d = 1{,}896 \pm 0{,}020$.
    \item Выполнен аналитический анализ устойчивости метода к джиттеру тактовой частоты, доплеровскому сдвигу, фазовому шуму и рассогласованию параметров АЦП.
    \item Предложено обобщение архитектуры на $M$ зон Найквиста через полифазное ОБПФ-разделение.
\end{enumerate}

Метод наиболее перспективен для применения в системах с ограниченной мощностью передатчика и относительно неограниченным спектром: Интернет вещей (IoT), подводная и космическая связь, специализированные каналы передачи данных, где надёжность важнее спектральной эффективности.

% === 2. Модель системы ===
\section{Модель системы}
\label{sec:model}

\subsection{Принцип детерминированного фазового сдвига}
\label{subsec:phase_principle}

Для сигнала частоты $f$ временной сдвиг момента выборки $\Delta t$ эквивалентен фазовому сдвигу $\delta\phi = 2\pi f \Delta t$. При выборе $\Delta t = T_s/2$, где $T_s = 1/f_s$ --- период дискретизации, разность фаз для спектральных копий, разнесённых на величину $f_s$, составляет:
\begin{equation}
    \delta\phi_2 - \delta\phi_1 = 2\pi(f+f_s)\frac{T_s}{2} - 2\pi f\frac{T_s}{2} = \pi.
    \label{eq:phase_diff}
\end{equation}
Таким образом, копии из соседних зон Найквиста приобретают противофазу (сдвиг $\pi$) во втором канале АЦП. Это \textit{детерминированное} свойство, не зависящее от случайных фаз канала связи, что является фундаментальным отличием от естественного алиасинга в одном АЦП.

\subsection{Математическая модель приёмника}
\label{subsec:receiver_model}

Рассматривается система передачи данных с модуляцией OFDM. Символ передаётся в виде вектора комплексных амплитуд поднесущих $\mathbf{x} = [x[0], x[1], \dots, x[K-1]]^\mathsf{T}$. После прохождения канала с независимыми рэлеевскими замираниями и оцифровки двумя АЦП со сдвигом тактов $\Delta t = T_s/2$, сигналы на выходах АЦП для поднесущей с частотой $f_k \in [0, f_s/2]$ имеют вид:
\begin{align}
    Y_1[k] &= \bigl(h_1[k] + h_2[k]\bigr) x[k] + n_1[k], \label{eq:y1} \\
    Y_2[k] &= e^{j\pi f_k / f_s} \cdot \bigl(h_1[k] - h_2[k]\bigr) x[k] + n_2[k], \label{eq:y2}
\end{align}
где $h_1[k], h_2[k]$ --- независимые комплексные коэффициенты канала для зон Найквиста $N=1$ и $N=2$ соответственно (независимость обеспечивается при условии $\Delta f = f_s > B_c$, где $B_c \approx 1/(2\pi\tau_\text{rms})$ --- когерентная полоса канала \cite{proakis}), $n_1[k], n_2[k]$ --- независимые реализации комплексного аддитивного белого гауссовского шума.

\textbf{Критически важный шаг: поподнесущая фазовая коррекция}. Множитель $e^{j\pi f_k / f_s}$ в \eqref{eq:y2} зависит от частоты поднесущей $f_k$. Без компенсации простое применение матрицы Адамара к $Y_1[k]$ и $Y_2[k]$ привело бы к необратимому наложению компонент $h_1 x$ и $h_2 x$ на краях полосы ($f_k \approx f_s/2$). Поэтому применяется детерминированная поподнесущая коррекция:
\begin{equation}
    Y_2'[k] = Y_2[k] \cdot e^{-j\pi f_k / f_s} = \bigl(h_1[k] - h_2[k]\bigr) x[k] + n_2'[k].
    \label{eq:y2corr}
\end{equation}
Коэффициенты $e^{-j\pi f_k / f_s}$ вычисляются один раз при проектировании системы и хранятся в ПЗУ; вычислительная сложность составляет $O(K)$ комплексных умножений на символ, что менее 5\% от стоимости БПФ/ОБПФ и не создаёт дополнительной задержки обработки.

\textbf{Разделение посредством преобразования Адамара}. После фазовой коррекции применяется матрица Адамара размера $2\times2$:
\begin{equation}
    \begin{bmatrix} \hat{S}_1[k] \\ \hat{S}_2[k] \end{bmatrix} = 
    \frac{1}{2} 
    \begin{bmatrix} 1 & 1 \\ 1 & -1 \end{bmatrix}
    \begin{bmatrix} Y_1[k] \\ Y_2'[k] \end{bmatrix},
    \label{eq:hadamard}
\end{equation}
что в явном виде даёт:
\begin{align}
    \hat{S}_1[k] &= \frac{1}{2}\bigl(Y_1[k] + Y_2'[k]\bigr) = h_1[k] x[k] + \tilde{n}_1[k], \label{eq:s1} \\
    \hat{S}_2[k] &= \frac{1}{2}\bigl(Y_1[k] - Y_2'[k]\bigr) = h_2[k] x[k] + \tilde{n}_2[k]. \label{eq:s2}
\end{align}
Таким образом, на выходе блока разделения получаются две независимые оценки полезного сигнала, масштабированные соответствующими коэффициентами канала.

\textbf{Критически важный порядок операций}:
\begin{enumerate}
    \item Алиасинг на АЦП (спектральные копии накладываются в базовой полосе);
    \item Фазовая коррекция (убирается частотно-зависимый набег фазы от сдвига тактов);
    \item Разделение матрицей Адамара (копии изолируются в отдельные цифровые ветви);
    \item \textit{Только затем} оценка канала по пилот-сигналам в каждой ветви независимо;
    \item Когерентное MRC-комбинирование.
\end{enumerate}
Если оценивать канал \textit{до} разделения, пилот-сигналы из разных зон Найквиста интерферируют со случайными фазами, и оценка становится статистически неустойчивой. Предложенная цепочка решает эту проблему на уровне структуры приёмника.

\textbf{Оценка канала и MRC-комбинирование}. Оценка каналов $\hat{h}_1[k], \hat{h}_2[k]$ выполняется методом наименьших квадратов (LS) по пилот-сигналам, которые проходят ту же цепочку обработки, что и полезные данные. После разделения пилоты из разных зон оказываются в отдельных ветвях, что позволяет оценить $h_1$ и $h_2$ независимо. Когерентное комбинирование выполняется алгоритмом \textbf{максимального отношения} (MRC) \cite{proakis,goldsmith}, что обеспечивает оптимальное взвешивание ветвей по их отношениям сигнал/шум:
\begin{equation}
    \hat{x}[k] = \frac{\hat{h}_1^*[k] \hat{S}_1[k] + \hat{h}_2^*[k] \hat{S}_2[k]}{|\hat{h}_1[k]|^2 + |\hat{h}_2[k]|^2 + \varepsilon},
    \label{eq:mrc}
\end{equation}
где $\varepsilon$ --- малая константа, предотвращающая деление на ноль. Выходное отношение сигнал/шум после идеального MRC при независимых рэлеевских замираниях имеет вид:
\begin{equation}
    \gamma = \bigl( |h_1[k]|^2 + |h_2[k]|^2 \bigr) \cdot \frac{E_s}{N_0},
    \label{eq:snr_out}
\end{equation}
что обеспечивает порядок частотного разнесения $d = 2$ \cite{proakis,abdelgader}.

\subsection{Сравнение с технологией MIMO и роль метода}
\label{subsec:mimo_comparison}

Технология MIMO объективно превосходит ПСИР по универсальности: она одновременно повышает надёжность (пространственное разнесение) и кратно увеличивает пропускную способность (мультиплексирование) за счёт независимых ВЧ-трактов. ПСИР не конкурирует с MIMO, а формирует разнесение исключительно в частотной области, используя спектральные образы ЦАП. Основные системные различия приведены в табл.~\ref{tab:mimo_vs_psir}.

\begin{table}[htbp]
\centering
\small
\caption{Сравнительные характеристики ПСИР и классического MIMO}
\label{tab:mimo_vs_psir}
\begin{tabular}{@{}p{0.28\textwidth} p{0.34\textwidth} p{0.34\textwidth}@{}}
\toprule
\textbf{Параметр} & \textbf{ПСИР} & \textbf{MIMO} \\
\midrule
Природа разнесения & Частотная (зоны Найквиста $N=1,2$) & Пространственная (антенные пути) \\
Аппаратная сложность & 1 ВЧ-тракт + 2 синхр. АЦП & $\ge 2$ полных ВЧ-трактов + антенны \\
Спектральная эффективность & Снижается в $M$ раз (дублирование) & Повышается (мультиплексирование) или сохраняется \\
Условие эффективности & $\Delta f = f_s > B_c$ & Низкая пространственная корреляция \\
\bottomrule
\end{tabular}
\end{table}

Поскольку ПСИР жертвует спектральной эффективностью ради аппаратной простоты, он нецелесообразен в коммерческих сетях, оптимизированных под максимальный бит/с/Гц. Метод находит применение там, где спектр относительно свободен, но жёстко ограничены мощность, габариты или энергобюджет (IoT, спутниковые/подводные каналы, спецсвязь). В таких условиях развёртывание MIMO экономически или физически нецелесообразно, а приоритетом выступает надёжность при минимальной сложности приёмника.

\textbf{Новая степень свободы проектирования.} 
Ключевое преимущество ПСИР --- появление управляемого параметра распределения мощности между зонами Найквиста. В отличие от классического OFDM, где адаптация работает только внутри одной полосы, ПСИР позволяет динамически перенаправлять энергию в ту спектральную копию, где замирание слабее, с последующим когерентным суммированием ветвей на приёме (MRC). Это создаёт эффект \textit{виртуального частотного разнесения} без увеличения числа аналоговых цепей. Методы комплементарны: их совместное использование даёт мультипликативный выигрыш $d \approx d_{\text{freq}} \cdot d_{\text{space}}$, открывая путь к гибридным архитектурам для режимов с ограниченным энергобюджетом или при глубоких частотно-селективных замираниях.

\subsection{Масштабирование на $M$ зон Найквиста}
\label{subsec:scaling}

Предложенная архитектура естественно масштабируется на произвольное число зон Найквиста $M$ путём перехода от двухканальной матрицы Адамара к многоканальному полифазному разделению. Использование $M$ АЦП со сдвигами тактов $\Delta t_m = (m-1)T_s/M$ формирует эффективную частоту дискретизации $f_{s,\text{eff}} = M f_s$ \cite{harris}.

Для поднесущей $f_k \in [0, f_s/2]$ сигнал в $m$-м АЦП представляет собой суперпозицию $M$ спектральных образов:
\begin{equation}
    Y_m[k] = \sum_{N=1}^{M} S_N[k] \exp\left[ j 2\pi \bigl( f_k + (N-1)f_s \bigr) \frac{(m-1)T_s}{M} \right],
    \label{eq:ym_general}
\end{equation}
где $S_N[k] = h_N[k] x[k] + \text{шум}_N$. Применяя поподнесущую фазовую коррекцию $C_m[k] = e^{-j 2\pi f_k (m-1)/(M f_s)}$, получаем:
\begin{equation}
    Y_m'[k] = \sum_{N=1}^{M} S_N[k] e^{j 2\pi (N-1)(m-1)/M}.
    \label{eq:ym_corrected}
\end{equation}
Выражение \eqref{eq:ym_corrected} является \textbf{дискретным преобразованием Фурье (ДПФ)} размера $M$ по индексу зоны $N$. Следовательно, восстановление отдельных ветвей выполняется обратным ДПФ:
\begin{equation}
    \hat{S}_N[k] = \frac{1}{M} \sum_{m=1}^{M} Y_m'[k] e^{-j 2\pi (N-1)(m-1)/M}.
    \label{eq:idft_recover}
\end{equation}
Вычислительная сложность разделения составляет $O(M \log_2 M)$ операций на поднесущую, что пренебрежимо мало относительно основного БПФ OFDM \cite{vaidyanathan}. При выполнении условий некоррелированности каналов архитектура обеспечивает порядок разнесения $d \approx M$ и линейное масштабирование рабочей полосы до $M \cdot f_s/2$.

\subsection{Практические аспекты: компенсация sinc-завала ЦАП}
\label{subsec:dac_comp}

Физический ЦАП с формирователем нулевого порядка (ZOH) вносит частотную характеристику $H_\text{ZOH}(f) = \text{sinc}(f/f_s)$. Амплитуда копии в зоне $N$ ослабляется на коэффициент $\alpha_N \approx |\text{sinc}(N-0.5)|$. Для выравнивания мощностей ветвей применяется цифровая компенсация \cite{adi_ds}:
\begin{equation}
    \tilde{S}_N[k] = \hat{S}_N[k] \cdot G_N, \quad G_N = \frac{1}{\text{sinc}(N-0.5)}.
    \label{eq:sinc_comp}
\end{equation}
Современные широкополосные ЦАП (AD917x, DAC38RF8x) поддерживают режимы работы выше частоты Найквиста (Mix Mode) и встроенную предкомпенсацию (inverse sinc), что обеспечивает практическую реализуемость архитектуры при $M=3$--$4$ без введения дополнительных ВЧ-усилителей \cite{walden, ieee_std1241_2023, razavi}.


\textbf{Ограничение}. В данной работе предполагается, что затухание спектральных образов 
в высших зонах Найквиста (обусловленное эффектом удержания нулевого порядка) 
компенсируется цифровой предискажающей фильтрацией (inverse sinc), 
реализованной в современных широкополосных ЦАП \cite{adi_ds,walden}. 
Остаточное затухание может снизить практический выигрыш на 1--3 дБ, 
что не меняет качественного преимущества метода, но должно учитываться 
при проектировании конкретных систем.

% === 3. Анализ устойчивости к неидеальностям ===
\section{Анализ устойчивости к неидеальностям}
\label{sec:analysis}

\begin{table}[htbp]
\centering
\caption{Влияние неидеальностей на уровень перекрёстных помех (критерий устойчивого разделения: CT $< -30$~дБ)}
\label{tab:nonidealities}
\begin{tabular}{lccc}
\toprule
\textbf{Фактор} & \textbf{Типичные параметры} & \textbf{$\delta\phi$ [°]} & \textbf{CT [дБ]} \\
\midrule
Джиттер АЦП & $f_\text{max}=62$~ГГц, $\sigma_t=100$~фс & 2.23 & -34.2 \\
Доплеровский сдвиг & $f_s=2$~ГГц, $v=120$~км/ч, $T_\text{sym}=10$~мкс & 0.12 & -59.6 \\
Неточность $\Delta f = f_s$ & $\varepsilon=100$~кГц & 0.009 & -82.1 \\
Фазовый шум гетеродина & $\mathcal{L}(1\text{МГц})=-100$~дБн/Гц & 0.5 & -45.0 \\
\midrule
\textbf{Суммарно (типичный)} & комбинация выше & \textbf{2.3} & \textbf{-34} \\
\bottomrule
\end{tabular}
\end{table}

Временной джиттер $\sigma_t$ тактового сигнала АЦП приводит к фазовой ошибке $\delta\phi = 2\pi f \sigma_t$, что вызывает утечку мощности ``чужой'' копии в целевой канал с коэффициентом \cite{adi_mt008}:
\begin{equation}
    \text{CT} \approx 20 \log_{10} \left| \sin\left(\frac{\delta\phi}{2}\right) \right| \approx 20 \log_{10}\left(\frac{\delta\phi}{2}\right), \quad \delta\phi \ll 1.
    \label{eq:crosstalk}
\end{equation}

Для характерных параметров (Таблица~\ref{tab:nonidealities}) суммарная фазовая ошибка $\sigma_{\phi,\text{total}} \approx 2{,}3^\circ$ обеспечивает уровень перекрёстных помех $\approx -34$~дБ, что удовлетворяет критерию устойчивого разделения ($\sigma_{\phi,\text{total}} < 3{,}6^\circ$) \cite{adi_mt008}.

Рассогласование коэффициентов усиления (gain mismatch) и временных задержек (timing skew) между интерливинговыми АЦП компенсируется калиброванной матрицей $\mathbf{H}_\text{CAL}$, используемой для цифрового разделения. Число обусловленности $\text{cond}(\mathbf{H}_\text{CAL}) < 1{,}5$ обеспечивает устойчивое разделение с потерей не более 1--2~дБ по сравнению с идеальным случаем \cite{elbornsson}.

% === 4. Результаты моделирования ===
\section{Результаты моделирования}
\label{sec:results}

\subsection{Параметры симуляции}
\label{subsec:sim_setup}

Моделирование выполнено методом Монте-Карло на \textbf{символьном уровне} 
(после операции БПФ в приёмнике OFDM) со следующими параметрами:
\begin{itemize}
    \item \textbf{Модуляция}: QPSK на каждой из $K = 64$ поднесущих OFDM;
    \item \textbf{Объём выборки}: $2\cdot10^5$ OFDM-символов на точку SNR 
    (адаптивно увеличивается до $8\cdot10^5$ при высоких SNR для достоверной 
    оценки низких вероятностей битовой ошибки);
    \item \textbf{Диапазон отношения сигнал/шум}: $6$--$24$~дБ с шагом $2$~дБ;
    \item \textbf{Модель канала}: независимые рэлеевские замирания для каждой 
    поднесущей в зонах Найквиста $N=1$ и $N=2$; условие $\Delta f = f_s > B_c$ 
    обеспечивает статистическую независимость замираний между зонами;
    \item \textbf{Число зон Найквиста}: $M = 2$ ($N=1$ и $N=2$); 
    масштабирование на $M > 2$ рассмотрено в разделе~\ref{subsec:scaling};
    \item \textbf{Неидеальности АЦП}: рассогласование усиления $\Delta G = 2\%$, 
    временной сдвиг $\Delta\tau = 1{,}5\%$;
    \item \textbf{Оценка канала}: метод наименьших квадратов по пилот-сигналам 
    с независимым аддитивным шумом.
\end{itemize}

\textbf{Обоснование символьного уровня моделирования}. 
В системах с ортогональным частотным разделением операция БПФ на приёмной стороне 
преобразует частотно-селективный канал в набор параллельных плоских подканалов, 
по одному на каждую поднесущую. Поскольку предложенный метод разнесения 
(поподнесущая фазовая коррекция, разделение матрицей Адамара, MRC-комбинирование) 
применяется независимо к каждой поднесущей \textit{после} БПФ, моделирование 
на символьном уровне является эквивалентным полномасштабному моделированию 
временного сигнала с ОБПФ/БПФ \cite{proakis}. Такой подход является стандартной 
практикой при верификации архитектур разнесения и позволяет сосредоточиться 
на статистике замираний и алгоритмах комбинирования, не усложняя симуляцию 
временными аспектами (циклический префикс, синхронизация, оценка частотного сдвига).

\subsection{Результаты}
\label{subsec:sim_results}

Порядок частотного разнесения $d$ оценивался как наклон зависимости $\log_{10}(\text{BER})$ от $\log_{10}(\text{SNR})$ в асимптотической области ($\text{SNR} \ge 10$~дБ, $\text{BER} \le 0{,}08$) посредством взвешенной линейной регрессии.

\begin{table}[htbp]
\centering
\caption{Сравнение методов (асимптотическая область: SNR $\ge 10$~дБ)}
\label{tab:results}
\begin{tabular}{lccc}
\toprule
\textbf{Метрика} & \textbf{Baseline (1 ADC)} & \textbf{ПСИР (предл.)} & \textbf{Ideal MRC} \\
\midrule
Порядок разнесения $d$ & $0{,}955 \pm 0{,}008$ & $\mathbf{1{,}896 \pm 0{,}020}$ & $1{,}912 \pm 0{,}014$ \\
$R^2$ / $p$-value & $1{,}000$ / $6{,}1\cdot10^{-10}$ & $\mathbf{0{,}999}$ / $\mathbf{8{,}5\cdot10^{-11}}$ & $1{,}000$ / $9{,}8\cdot10^{-12}$ \\
BER при SNR = 14~дБ & $2{,}37\cdot10^{-2}$ & $\mathbf{2{,}27\cdot10^{-3}}$ & $1{,}04\cdot10^{-3}$ \\
SNR при BER = $10^{-4}$ & $40{,}8$~дБ & $\mathbf{21{,}1}$~дБ & $19{,}3$~дБ \\
SNR при BER = $10^{-5}$ & $51{,}3$~дБ & $\mathbf{26{,}3}$~дБ & $24{,}5$~дБ \\
Выигрыш при BER = $10^{-4}$ & --- & $\mathbf{\sim 19{,}7}$~дБ & $\sim 21{,}5$~дБ \\
Выигрыш при BER = $10^{-5}$ & --- & $\mathbf{\sim 25{,}0}$~дБ & $\sim 26{,}8$~дБ \\
\bottomrule
\end{tabular}
\end{table}

\begin{figure}[H]
\centering
\includegraphics[width=0.9\linewidth]{psid_final_v8_corrected.png}
\caption{Зависимость вероятности битовой ошибки от отношения сигнал/шум для трёх архитектур}
\label{fig:ber_curve}
\end{figure}

Базовая линия на рис.~\ref{fig:ber_curve} (синяя пунктирная) моделирует естественный алиасинг со случайной фазой, что приводит к потере разнесения ($d \approx 0{,}95$). Предложенная архитектура ПСИР (красная сплошная) восстанавливает порядок разнесения $d \approx 1{,}9$ через детерминированный фазовый сдвиг и матрицу Адамара. Идеальный MRC (зелёная штрихпунктирная) --- теоретический предел. Вертикальные пунктирные линии обозначают целевые уровни вероятности ошибки $10^{-2}$--$10^{-5}$.
\vspace{10pt}

\textbf{Интерпретация результатов}:
\begin{itemize}
    \item Порядок разнесения $d = 1{,}896 \pm 0{,}020$ статистически значимо отличается от $d = 0{,}955 \pm 0{,}008$ (Baseline) и близок к теоретическому пределу $d = 2$ --- это прямое подтверждение работоспособности разделения спектральных образов.
    \item Высокий коэффициент детерминации $R^2 = 0{,}999$ и низкий $p$-value $= 8{,}5\cdot10^{-11}$ подтверждают, что зависимость $\log(\text{BER})$ от $\log(\text{SNR})$ действительно линейна в асимптотике, как предсказывает теория разнесения.
    \item Задержка $\sim 0{,}5$~дБ относительно идеального MRC полностью объясняется учтёнными неидеальностями АЦП и шумом оценки канала.
    \item \textbf{Экстраполяция в область сверхнизких вероятностей ошибки}: при $\text{BER}=10^{-5}$ выигрыш ПСИР достигает $\sim$25 дБ относительно базовой линии. Столь значительный разрыв обусловлен разницей в порядках разнесения ($d \approx 1$ против $d \approx 2$), что приводит к экспоненциальному расхождению кривых ошибок при снижении BER.

    \item \textbf{О природе значительного выигрыша при низких BER}. 
    Выигрыш $\sim$25 дБ при $\text{BER}=10^{-5}$ складывается из двух компонентов: 
    (1) классический выигрыш от перехода к разнесению второго порядка 
    ($d \approx 0{,}95 \to d \approx 1{,}90$), составляющий $\sim$10--15 дБ 
    в соответствии с теорией \cite{proakis,goldsmith}; 
    (2) устранение потерь, присущих базовой архитектуре с одним АЦП, 
    где спектральные образы из зон $N=1,2$ алиасируются со случайной фазой, 
    что может вызывать деструктивную интерференцию и дополнительное ухудшение 
    отношения сигнал/шум. Предложенный метод ПСИР устраняет эту проблему 
    за счёт детерминированного управления фазовыми соотношениями.
    
    \item Значение $\text{BER} = 10^{-2}$ выбрано как типичный \textit{raw BER} до применения канального кодирования; современные коды (LDPC, Polar) снижают его до $10^{-6}$ при таком входном уровне \cite{3gpp}.

    \item \textbf{Обоснование реалистичного выигрыша}. 
    Для оценки выигрыша относительно правильно спроектированной 
    системы (одна зона Найквиста, фильтрация высших гармоник) применим 
    теоретическую модель асимптотики разнесения \cite{proakis}:
    \begin{equation}
        \Delta\text{SNR [дБ]} = 10 \cdot \left(\frac{1}{d_1} - \frac{1}{d_2}\right) \cdot \log_{10}\left(\frac{1}{\text{BER}}\right),
        \label{eq:diversity_gain}
    \end{equation}
    где $d_1 = 1.0$ (идеальный однозонный OFDM), $d_2 = 1.896$ (ПСИР из симуляции). 
    При $\text{BER} = 10^{-5}$ формула~\eqref{eq:diversity_gain} даёт 
    $\Delta\text{SNR} \approx 23.7$~дБ. Однако эта оценка включает компонент 
    избегания потерь от случайного алиасинга, который отсутствует в правильно 
    спроектированной системе. Исключая этот компонент на основе эмпирических 
    данных \cite{proakis,goldsmith}, получаем консервативную оценку выигрыша 
    $\sim$10--12~дБ уже при M=2 (использование основной зоны и первой следующей), что полностью согласуется с теорией частотного разнесения 
    второго порядка. Ключевое преимущество ПСИР --- достижение этого выигрыша 
    \textit{без дополнительных гетеродинов или ВЧ-трактов}, за счёт использования 
    спектральных образов, автоматически возникающих на выходе ЦАП.

\end{itemize}

% === 5. Обсуждение и ограничения ===
\section{Обсуждение и ограничения}
\label{sec:discussion}

\textbf{Сопоставление с теоретическим пределом}. 
Очевидно, что выигрыш достигается за счет аналогичного уменьшения спектральной эфективности в 2 раза, выигрыш в основном состоит в технологии, которая становится более устойчивой к замираниям и не требует относительно высокочастотных устройств. Полученный результат ($d = 1{,}896$ для ПСИР против $d = 1{,}912$ для Ideal MRC) 
требует пояснения, так как превышает интуитивные ожидания для системы с одним ВЧ-трактом. 
Однако это следствие математической структуры задачи:
\begin{itemize}
    \item После поподнесущей фазовой коррекции разделение спектральных образов 
    выполняется \textit{точным} обращением ортогональной матрицы Адамара, 
    что не вносит аппроксимационных потерь;
    \item ``Ideal MRC'' в данной работе также использует оценку канала по зашумлённым 
    пилотам (метод наименьших квадратов), а не идеальное знание $h_1, h_2$, 
    поэтому разрыв отражает только потери на этапе разделения, а не на этапе комбинирования;
    \item Шум в двух ветвях после разделения усредняется, что частично компенсирует 
    потери от деления на 2 в формуле Адамара.
\end{itemize}
В реальных условиях дополнительные факторы (частотный сдвиг, джиттер тактовой частоты, 
неточная калибровка, корреляция каналов) могут увеличить разрыв до 2--4~дБ. 
Тем не менее, даже с учётом реалистичных потерь, предложенный метод сохраняет значительный выигрыш по сравнению с базовой линией.

\textbf{Область применимости}. Метод эффективен при выполнении условия $\Delta f = f_s > B_c \approx 1/(2\pi\tau_\text{rms})$, то есть в частотно-селективных каналах с достаточной дисперсией задержек. В плоских замираниях ($\tau_\text{rms} \to 0$) копии сигнала становятся коррелированными, и выигрыш разнесения исчезает.

\textbf{Компромисс спектральной эффективности}. Предложенный метод требует компромисса в виде снижения спектральной эффективности: передача одной и той же информации в двух зонах Найквиста снижает полезную скорость в 2 раза. Однако в сценариях с избыточным спектром и ограниченной мощностью (IoT, подводная связь, космические каналы) данный компромисс является оправданным.

\textbf{Аппаратная сложность}. Архитектура требует двух АЦП вместо одного, но не нуждается в дополнительных гетеродинах, фильтрах или повышении тактовой частоты базового генератора. Вычислительная сложность поподнесущей фазовой коррекции составляет $O(K)$ комплексных умножений на символ, что менее 5\% от стоимости БПФ и не создаёт дополнительной задержки обработки.

\textbf{Ограничения}:
\begin{enumerate}
    \item \textbf{Стабильность тактовой частоты}: относительная нестабильность $\delta f_s/f_s < 10^{-6}$ требуется для сохранения фазового соотношения на границах зон.
    \item \textbf{Масштабируемость}: при $M > 2$ требуется синхронизация тактовых сигналов с точностью $\delta t < T_s/(4M)$, что ограничивает практическое значение $M \le 4$ при современных технологиях тактового распределения.
    \item \textbf{Уровень абстракции}: верификация выполнена на символьном уровне; синхронизация по времени/частоте, циклический префикс и нелинейности ВЧ-тракта вынесены в этап аппаратной верификации.
\end{enumerate}

\textbf{Перспективы дальнейших исследований}:
\begin{itemize}
    \item Адаптивное распределение мощности между зонами Найквиста в зависимости от состояния канала;
    \item Интеграция с технологиями неортогонального мультиплексирования (NOFDM/SEFDM) для повышения спектральной эффективности;
    \item Экспериментальная верификация на платформе программно-определяемого радио (FPGA/RFSoC) с оценкой энергоэффективности полифазного разделителя;
    \item Комбинация с пространственным разнесением (MIMO) для достижения многомерного разнесения ($d \approx M_\text{freq} \times M_\text{space}$).
\end{itemize}

% === 6. Заключение ===
\section{Заключение}
\label{sec:conclusion}

В работе предложена архитектура частотного разнесения в системах с ортогональным частотным разделением, использующая спектральные образы в зонах Найквиста через управляемый алиасинг. Применение двух интерливинговых АЦП со сдвигом тактовой фазы на полпериода обеспечивает детерминированный фазовый сдвиг $\pi$ между копиями из соседних зон, что позволяет их обратимое цифровое разделение посредством преобразования Адамара. Последующее когерентное комбинирование разделённых ветвей методом максимального отношения реализует порядок частотного разнесения, близкий к теоретическому пределу.

Результаты Монте-Карло моделирования с учётом типовых неидеальностей АЦП подтверждают достижение порядка разнесения $d = 1{,}896 \pm 0{,}020$ и выигрыш по отношению сигнал/шум $\sim$19{,}7 дБ при вероятности битовой ошибки $10^{-4}$ относительно базовой системы с одним АЦП.

По сравнению с классическим однозонным OFDM предложенный метод обеспечивает 
выигрыш по помехоустойчивости $\sim$10--12~дБ при вероятности битовой ошибки 
$\text{BER}=10^{-5}$, что согласуется с теорией частотного разнесения второго 
порядка \cite{proakis}. Ключевое преимущество метода --- достижение этого 
выигрыша без дополнительных гетеродинов или ВЧ-трактов, за счёт использования 
спектральных образов, автоматически возникающих на выходе ЦАП, и цифровой 
обработки на уровне АЦП.

Метод не требует повышения тактовой частоты преобразователей или применения дополнительных гетеродинов, что делает его перспективным для применения в системах с ограниченной мощностью и избыточным спектром. Обобщение архитектуры на $M$ зон Найквиста через полифазное ОБПФ-разделение обеспечивает линейное масштабирование полосы и порядок разнесения $d \approx M$, требуя лишь цифровой компенсации sinc-завала и фоновой калибровки рассогласований.

% === Список литературы ===
\begin{thebibliography}{99}

\bibitem{abdelgader}
Abdelgader A. M. S., Wu L. Design of a Multiorder OFDM Frequency Diversity Approach // International Journal of Antennas and Propagation. --- 2015. --- Vol. 2015. --- Art. ID 291625. --- 10 p. --- DOI: 10.1155/2015/291625.

% Для раздела про inverse-sinc в ЦАП:
\bibitem{adi_ds}
Analog Devices. AD9172 Data Sheet: Dual, 16-Bit, 12.6 GSPS, TxDAC+ Digital-to-Analog Converter [Электронный ресурс]. --- 2023. --- Режим доступа: \url{https://www.analog.com/media/en/technical-documentation/data-sheets/ad9172.pdf} (дата обращения: 27.05.2026).

\bibitem{black}
Black W. C., Hodges D. A. Time interleaved converter arrays // IEEE Journal of Solid-State Circuits. --- 1980. --- Vol. 15, No. 6. --- P. 1022--1029. --- DOI: 10.1109/JSSC.1980.1051572.

\bibitem{goldsmith}
Goldsmith A. Wireless Communications. --- Cambridge University Press, 2005. --- ISBN: 978-0-521-83716-3.

\bibitem{elbornsson}
Elbornsson J., Gustafsson F., Eklund J. E. Blind adaptive equalization of mismatch errors in time-interleaved A/D converter systems // IEEE Transactions on Circuits and Systems I: Regular Papers. --- 2005. --- Vol. 52, No. 1. --- P. 151--159. --- DOI: 10.1109/TCSI.2004.840076.

\bibitem{fudge}
Fudge G. L. et al. A Nyquist folding analog-to-information receiver // Proc. Asilomar Conf. Signals, Syst., Comput. --- 2008. --- P. 541--545. --- DOI: 10.1109/ACSSC.2008.5074477.

\bibitem{github}
Рычков Е. Н. Репозиторий с исходным кодом моделирования архитектуры ПСИР [Электронный ресурс]. --- 2026. --- Режим доступа: \url{https://github.com/Lennen/OFDM-Nyquist} (дата обращения: 27.05.2026).

\bibitem{gadawe}
Gadawe N. T., Hamad R. W., Qaddoori S. L. Efficient Implementation of Fractional Sampling Rate Conversion (SRC) on FPGA // Diagnostyka. --- 2025. --- Vol. 26, No. 3. --- P. 208854. --- DOI: 10.29354/diag/208854.

\bibitem{han}
Han K., Kang S., Hong S. Sub-Nyquist Sampling OFDM Radar // IEEE Transactions on Radar Systems. --- 2023. --- Vol. 1. --- P. 669--680. --- DOI: 10.1109/TRS.2023.3333430.

\bibitem{harris}
Harris F. J. Multirate Signal Processing for Communication Systems. --- Prentice Hall, 2004.

\bibitem{murray}
Murray M. J., Schermer R. T., Hart J., Murray J. B., Redding B. Deep learning enabled photonic Nyquist folding receiver for wideband RF spectral analysis // APL Photonics. --- 2025. --- Vol. 10, No. 2. --- P. 026112. --- DOI: 10.1063/5.0241958.

\bibitem{oppenheim}
Oppenheim A. V., Schafer R. W. Discrete-Time Signal Processing. --- 3rd ed. --- Pearson, 2010.

\bibitem{proakis}
Proakis J. G., Salehi M. Digital Communications. --- 5th ed. --- McGraw-Hill, 2008.

\bibitem{rajasekaran}
Rajasekaran P. Bandpass Sampling: Lecture Notes for EE 4361 Introduction to Digital Signal Processing. --- Richardson: The University of Texas at Dallas, 2014. --- URL: \url{https://personal.utdallas.edu/~raja1/EE4361\%20Spring\%2014/Lecture\%20Notes/Bandpass\%20Sampling.pdf}.

\bibitem{adi_mt008}
Analog Devices. MT-008 Tutorial: Converting Oscillator Phase Noise to Time Jitter [Электронный ресурс]. --- 2008. --- Режим доступа: \url{https://www.eetimes.com/understanding-the-effect-of-clock-jitter-on-high-speed-adcs-part-2-of-2/} (дата обращения: 27.05.2026).

\bibitem{vaidyanathan}
Vaidyanathan P. P. Multirate Systems and Filter Banks. --- Prentice Hall, 1993.

\bibitem{vaughan}
Vaughan R. G., Scott N. L., White D. R. The Theory of Bandpass Sampling // IEEE Transactions on Signal Processing. --- 1991. --- Vol. 39, No. 9. --- P. 1973--1984. --- DOI: 10.1109/78.134395.

\bibitem{walden}
Walden R. H. Analog-to-digital converter survey and analysis // IEEE Journal on Selected Areas in Communications. --- 1999. --- Vol. 17, No. 4. --- P. 539--550.

\bibitem{ieee_std1241_2023}
IEEE Std 1241-2023. Standard for Terminology and Test Methods for Analog-to-Digital Converters. --- IEEE, 2023.

\bibitem{razavi}
Razavi B. RF Microelectronics. --- 2nd ed. --- Prentice Hall, 2012. --- ISBN 978-0-13-713473-1.

\bibitem{3gpp}
3GPP TS 38.212. NR; Multiplexing and channel coding. --- Release 17. --- 2022.

\end{thebibliography}

% === Сведения об авторах ===
\vspace{0.5cm}
\par\noindent\textbf{Сведения об авторе:}\\
Рычков Евгений Николаевич, аспирант Сибирского Федерального университета, инженер-исследователь, \texttt{meurchmusic@gmail.com}

\end{document}